\subsubsection{BlackBerry OS}

\paragraph{File System Structure and Organization}

BlackBerry~10/QNX's flagship file system is the Power-Safe file system
(\texttt{fs-qnx6}), a reliable on-disk file system explicitly designed
to survive power failures without losing or corrupting data
\citep{qnx_powersafe_fs}. BlackBerry~10 is also POSIX-based, so it provides transparent access to
other file systems, including DOS/FAT and Linux ext2
\citep{qnx_neutrino_overview, qnx_fs_fault_tolerance}. Power-Safe is inode-based like traditional UNIX file
systems and is created with \texttt{mkqnx6fs}, which lets you set
logical block size, endianness, block count, and the maximum number of
inodes (hence files) \citep{qnx_powersafe_fs}. Its defining structural
property is copy-on-write (COW): modified blocks are relocated rather
than overwritten in place \citep{qnx_fs_fault_tolerance}. The advantage
is integrity; the limitation is the write amplification and bookkeeping
COW entails.

\paragraph{File Access Methods and Operations}

Basic operations (\texttt{open}, \texttt{read}, \texttt{write},
\texttt{close}) follow the POSIX model, but the mechanism is pure
microkernel: a file operation is a message to a resource-manager
process (the file system running in user space), not a kernel call
\citep{qnx_memman, qnx_powersafe_fs}. The OS supports sequential and
random access, and processes hold file descriptors that map to
connections to the file-system resource manager. Internally, when a
file is accessed the request is sent to \texttt{fs-qnx6}, which
consults its cached metadata and buffer cache, performs the read/write
against relocated COW blocks, and replies---the same Send/Receive/Reply
pattern that governs all QNX IPC \citep{qnx_ipc, qnx_powersafe_fs}.

\paragraph{Directory Structures}

BlackBerry~10 uses a hierarchical, tree-structured namespace rooted at
\texttt{/}, with \texttt{/} (0x2F) as the path separator
\citep{qnx_powersafe_fs}.
Directories are implemented as special files containing name-to-inode
entries; path resolution walks the tree component by component, and
because the namespace is unified, a path may even resolve across the
network to another node transparently \citep{qnx_ipc, qnx_powersafe_fs}.
The performance implication is that deep paths cost more lookups,
mitigated by metadata caching in the file-system process.

\paragraph{Disk Space Allocation and Free Space Management}

Power-Safe uses indexed (inode/extent-style) allocation and tracks free
space with on-disk structures updated through COW. Crucially, space is
allocated by relocating modified blocks to free space and only then
committing a new superblock/root that references them, so at any
instant either the complete old version or the complete new version is
valid---never a half-written mix \citep{qnx_fs_fault_tolerance}. File
growth and shrinking are handled by allocating/relocating blocks;
Power-Safe also creates sparse files by default when a file is grown
with \texttt{ftruncate()} \citep{qnx_powersafe_fs}. The trade-off
versus simple in-place allocation is extra space and write activity in
exchange for crash safety.

\paragraph{File Protection and Security}

BlackBerry~10 enforces POSIX/UNIX-style permissions (owner/group/other
read--write--execute) plus user and group ownership, backed by the
user/kernel-mode and MMU isolation already described
\citep{qnx_memman, qnx_neutrino_overview}. The Power-Safe file system
additionally offers encryption of all or part of the file system
\citep{qnx_powersafe_fs}. These mechanisms give a familiar,
well-understood security model; their limitation is that classic UNIX
permissions are coarser than full access-control lists (ACLs), though
they are simple and fast.

\paragraph{Reliability and Crash Recovery}

This is QNX's standout file-system strength. Rather than journaling,
Power-Safe uses copy-on-write plus periodic snapshots: a snapshot
captures a consistent view of the whole file system, taken
automatically (every 10 seconds by default), and also on
\texttt{sync()} or \texttt{fsync()}
\citep{qnx_powersafe_fs, qnx_fs_fault_tolerance}. If power is lost
mid-write, the half-written new blocks are simply never committed, and
the system falls back to the last consistent snapshot---``at any point
in time either the old version is completely accessible or the new
version is completely accessible,'' with no live data overwritten in
between \citep{qnx_fs_fault_tolerance}. This specifically defends
against the corrupted ECC sectors that ordinary file systems can
suffer on power loss \citep{qnx_powersafe_fs}. The trade-off:
durability is bounded by snapshot frequency (work since the last
snapshot can be lost), and synchronous guarantees require drives that
honour cache-flush/synchronisation commands
\citep{qnx_fs_fault_tolerance}.

\paragraph{Disk Scheduling and I/O Management}

Disk I/O is mediated by the block I/O driver stack (\texttt{devb-*})
and the file-system resource manager, which queue and order requests
to the underlying device \citep{qnx_powersafe_fs}. Buffering and a
configurable buffer cache sit between applications and the disk, and
snapshot/commit timing governs when dirty blocks are flushed
\citep{qnx_powersafe_fs, qnx_powersafe_perf}. Because QNX is real-time
oriented, its I/O path emphasises predictable latency and integrity
over squeezing out maximum raw throughput.

\paragraph{Performance Considerations}

Performance ``depends on how much buffer cache is available, and on
the frequency of the snapshots'' \citep{qnx_powersafe_perf}. More cache
and less frequent snapshots raise throughput but widen the window of
potential data loss and the cost of an eventual sync; conversely,
frequent snapshots maximise safety at a throughput cost---and because
synchronisation is at the \emph{file-system} level, \texttt{fsync()} is
potentially expensive and \texttt{O\_SYNC} is ignored
\citep{qnx_powersafe_perf}. Read-ahead and buffering help large
sequential files; many small files stress metadata and inode handling.
The headline trade-off is that crash safety (COW + snapshots) costs
some write performance.

\paragraph{Overview --- File and Disk Management}

Strengths include outstanding crash/power-failure resilience via COW
and snapshots, integrity by design, POSIX semantics, multi-file-system
interoperability (ext2/FAT/exFAT), and optional encryption
\citep{qnx_powersafe_fs, qnx_neutrino_overview, qnx_fs_fault_tolerance}.
Weaknesses are COW write amplification and snapshot-bounded durability;
a coarse permission model relative to full ACLs; and tuning for
reliability over peak speed
\citep{qnx_fs_fault_tolerance, qnx_powersafe_perf}. For a logging
system doing frequent writes, Power-Safe protects integrity but the
constant block relocation and periodic full-filesystem sync impose
overhead---a designer would tune snapshot frequency and cache to
balance safety against write rate \citep{qnx_powersafe_perf}. For large
multimedia files, sequential read-ahead and buffering serve well,
though a throughput-optimised competitor might edge ahead on pure
streaming speed where crash safety is less critical.
